\subsection{バスの時刻と組み合わせた提示する時刻の取得}
% +10してその帰宅推奨時刻に帰ったユーザが乗ったと思われるバスの時刻に調整した話入れないといけない，根拠も
% 二択の場合許容，並び替え時点で最後の選択消したね
% 帰宅推奨時刻に最も近いバスの時刻とその前後のバスの時刻を取得する．
% 以降の処理では，上記の条件を満たし，有効な帰宅推奨時刻が
% 得られた場合のみ，バス時刻との組み合わせによる
% 提示時刻の算出を行う．
% 選択肢を用意し，よりユーザの都合に合わせやすい設計にした
% APIで取得できたらいいよね，実際にこんなAPIあるよねという引用もしたい
算出された帰宅推奨時刻に対して，その時刻に最も近いバスの発車時刻を基準とし，さらにその前後に位置するバスの発車時刻を含めた三つの時刻を取得する．
前節までの手法により有効な帰宅推奨時刻が算出できた場合に限り，バスの時刻との組み合わせによる提示時刻の算出を行う．
なお，本研究ではバス時刻情報を外部から取得する前提で設計しており，実運用においては公共交通機関が提供する時刻表API等との連携を想定している．
複数の選択肢の提示によりユーザごとの準備状況や行動の柔軟性に対応し，帰宅行動の受容性の向上を目的としている．
また，提示する時刻は時間順に整列した上で表示する．

帰宅推奨時刻とバス発車時刻との対応付けにあたっては，ユーザの移動時間を考慮した調整を行っている．
具体的には，帰宅推奨時刻に一定時間を加えた時刻を基準として，最も近いバスの発車時刻を取得する．
この加算時間は利用環境に依存するパラメータであり，本研究では対象環境における実測値に基づいて設定している．

提示可能なバス時刻の組合せは，帰宅推奨時刻に対応するバスの運行状況により，三択とならない場合がある．
帰宅推奨時刻の直後に運行されるバスが存在しない場合には，直前に位置するバスの時刻を含めた二つの時刻を提示する．
二択提示は，ユーザに一定の時間的制約を意識させつつも，完全に行動を固定しない選択の余地を残すため，帰宅行動を促進する効果が期待される．
そこで，本研究では有効な提示形式として採用する．
一方，提示可能な時刻が一つのみとなる場合には，選択の余地がないと判断し，提示対象から除外する．

なお，本システムでは，提示されるバス時刻のうち最も早い時刻に対して，投票結果を確認した上で退室しても当該バスに間に合うよう，一定の余裕時間を設けた投票締切時刻を設定している．
この締切時刻は，ユーザが意思決定を行うための最終時刻として機能し，帰宅行動の現実的な成立を目的としている．
この締切時刻も同様に利用環境に依存するパラメータであり，本研究では対象環境における実測値に基づいて設定している．